\subsection{Any number of rows with two minimal columns}

\begin{theorem}[Proved: an identity-adjoined two-column band opposite a graph corner]
\label{CP2:A-two-column-band-graph-source}
Let $M=K^1$, where
$K=\{e_{ij}:0\le i<r,\ j\in\{0,1\}\}$ is a rectangular
band, $r\ge1$, with products $e_{ij}e_{kl}=e_{il}$.
Suppose a finite nonempty left-$N$, right-$M$ biset has
the natural equivariant graph comparison and rationally
acyclic opposite corner used in
Theorem~\ref{CP2:A-projective-local-row-graph-obstruction}.
Then all rational degree-two products of the collage vanish.
Boundedness of its homology is not required.

In particular this holds for $N=R_3$ and every finite
connecting biset.  For $r=3$, it excludes source $42$
of the frozen positive order-seven catalog from the
eighteen-arrow budget $(7,3,8)$.
\end{theorem}
\begin{proof}
Let $R=\mathbb QM$, put $e=e_{00}$ and $e'=e_{01}$,
and use the natural augmented graph-chain comparison $\Gamma$
for the opposite-corner action.  Collapsing the rationally
acyclic opposite corner gives an actual rational homology
equivalence with
\[
 Y=\operatorname{hocofib}(B(N,S,M)\longrightarrow BM).
\]
This actual comparison induces a rational cohomology-ring
isomorphism.  The graph comparison is used to compute
augmented rational chains of this actual cofiber; it need
not itself be a comparison of spaces.
As in the cited graph calculation, write
\[
 C=[C_1(\Gamma;\mathbb Q)\longrightarrow
                  \ker(C_0(\Gamma;\mathbb Q)\to\mathbb Q)]
\]
in degrees one and zero.  The reduced chains of $Y$ are
computed by $(C\otimes_R^{\mathbf L}\mathbb Q)[1]$.

If $\Gamma e$ and $\Gamma e'$ have a common vertex $v$,
then $ve=v=ve'$.  For every band element,
\[
 ve_{ij}=vee_{ij}=ve_{0j}=v.
\]
Thus $v$ is fixed by all of $M$.  Its augmentation section,
transported along the natural graph-chain comparison,
makes the actual bridge pullback injective on rational
$H^2(BM)$.  Every second class of $Y$ is then relative
to $BM$, and the actual suspension quotient kills all
their products.  This argument uses no homological bound.

Suppose now that the two image graphs have disjoint
vertex sets, and therefore disjoint edge sets.  The
image of any band actor with column $j$ is the same
as the image of $e_{0j}$: the identities
\[
 e_{ij}e_{0j}=e_{ij},\qquad e_{0j}e_{ij}=e_{0j}
\]
give the two image inclusions.  Consequently multiplication
gives a degreewise injective chain map
\[
 J=Ce\otimes_{\mathbb Q}eR\longrightarrow C
\]
with image $C\mathbb QK$.  Indeed $eR$ has the two
column basis elements $e,e'$, and the corresponding
copies of $Ce$ have disjoint supports.  In degree zero
the vectors in both copies have augmentation zero;
in degree one disjoint supports also prevent cancellation.
Every term of $J$ is a projective right $R$ module.
The quotient $D=C/J$ is a two-term right
\[
 A=R/\mathbb QK\cong\mathbb Q
\]
complex.  Removing $J$ changes no derived tensor homology
in degrees at least three.

The left projective $Re$ has basis $e_{i0}$, $0\le i<r$.
Its augmentation gives the exact sequence
\[
 0\longrightarrow V\longrightarrow Re
   \longrightarrow\mathbb Q\longrightarrow0,
 \qquad
 V=\mathbb Q\{e_{i0}-e_{00}:1\le i<r\}.
\]
Every band element annihilates $V$, so it is the free
left $A=\mathbb Q$ module of dimension $r-1$.
Since $De=0$, this sequence identifies
$D\otimes_R^{\mathbf L}\mathbb Q$ with
$(D\otimes_R^{\mathbf L}V)[1]$.
For completeness, the ideal has the factorization
\[
 \mathbb QK\cong Re\otimes_{\mathbb Q}eR
\]
and is projective as a left $R$ module.
Thus $0\to\mathbb QK\to R\to A\to0$ is a
left-projective resolution of $A$, and
$D\otimes_R^{\mathbf L}\mathbb QK=0$ because $De=0$.
It follows that $D\otimes_R^{\mathbf L}A\simeq D$,
and derived associativity identifies the preceding
tensor with $D\otimes_A^{\mathbf L}V$.

Combining the two shifts gives the natural equality
\[
 H_n(Y;\mathbb Q)
   =H_{n-2}(D\otimes_A^{\mathbf L}V)
                         \qquad(n\ge4).
\]
The right side is zero because $V$ is projective over
$A$ and $D$ has degrees zero and one.  Hence
$H_4(Y;\mathbb Q)=0$, which kills all degree-two
products in this second case as well.  The row count
never enters the column-injectivity argument; it only
changes the free multiplicity of $V$.

For $N=R_3$, the required natural graph comparison is
the one used in
Corollary~\ref{CP2:A-baseline-right-zero-three-six};
its construction depends on the opposite $R_3$ action,
not on the order of $M$.
The source-$42$ identification has identity $0$,
$r=3$, and band labels $e_{ij}=1+2i+j$.
All forty-nine products are reconstructed in
\path{converse/scope18/two_column_band_source42_literal_binding_v2.json}.
\end{proof}
